% !TeX root = ../sections/03_solutions.tex
\subsection{1.1.A.6.2}
\begin{proof}
(1) より，ある自然数 $N_0$ が存在して，$n\geq N_0$ ならば
    \[
        |b_n|\geq \frac{|\beta|}{2}
    \]
    が成り立つ。

    示したいことは
    \[
        \lim_{n\to\infty}\frac{a_n}{b_n}
        =
        \frac{\alpha}{\beta}
    \]
    である。すなわち，任意の $\epsilon>0$ に対して，ある自然数 $N$ が存在して，
    $n\geq N$ ならば
    \[
        \left|
            \frac{a_n}{b_n}-\frac{\alpha}{\beta}
        \right|<\epsilon
    \]
    となることを示せばよい。

    まず差を通分する。
    \[
        \frac{a_n}{b_n}-\frac{\alpha}{\beta}
        =
        \frac{\beta a_n-\alpha b_n}{\beta b_n}
    \]

    分子を次のように変形する。
    \[
        \beta a_n-\alpha b_n
        =
        \beta a_n-\beta\alpha+\beta\alpha-\alpha b_n
    \]
    したがって，
    \[
        \beta a_n-\alpha b_n
        =
        \beta(a_n-\alpha)+\alpha(\beta-b_n)
    \]
    である。

    よって，
    \[
        \left|
            \frac{a_n}{b_n}-\frac{\alpha}{\beta}
        \right|
        =
        \left|
            \frac{\beta(a_n-\alpha)+\alpha(\beta-b_n)}{\beta b_n}
        \right|
    \]
    である。
    三角不等式を用いると，
    \[
        \left|
            \frac{a_n}{b_n}-\frac{\alpha}{\beta}
        \right|
        \leq
        \frac{|\beta||a_n-\alpha|+|\alpha||\beta-b_n|}
        {|\beta||b_n|}
    \]
    となる。

    ここで，$n\geq N_0$ ならば
    \[
        |b_n|\geq \frac{|\beta|}{2}
    \]
    であるから，
    \[
        \frac{1}{|b_n|}
        \leq
        \frac{2}{|\beta|}
    \]
    が成り立つ。

    したがって，$n\geq N_0$ のとき
    \[
        \left|
            \frac{a_n}{b_n}-\frac{\alpha}{\beta}
        \right|
        \leq
        \frac{2}{|\beta|}|a_n-\alpha|
        +
        \frac{2|\alpha|}{|\beta|^2}|b_n-\beta|
    \]
    を得る。

    ここで，任意に $\epsilon>0$ を取る。

    $a_n\to\alpha$ であるから，ある自然数 $N_1$ が存在して，
    $n\geq N_1$ ならば
    \[
        |a_n-\alpha|<\frac{|\beta|\epsilon}{4}
    \]
    が成り立つ。

    また，$b_n\to\beta$ であるから，ある自然数 $N_2$ が存在して，
    $n\geq N_2$ ならば
    \[
        |b_n-\beta|
        <
        \frac{|\beta|^2\epsilon}{4(|\alpha|+1)}
    \]
    が成り立つ。

    ここで
    \[
        N=\max\{N_0,N_1,N_2\}
    \]
    とおく。

    このとき，$n\geq N$ ならば，
    \[
        n\geq N_0,\qquad n\geq N_1,\qquad n\geq N_2
    \]
    である。

    したがって，
    \[
        \frac{2}{|\beta|}|a_n-\alpha|
        <
        \frac{2}{|\beta|}\cdot \frac{|\beta|\epsilon}{4}
        =
        \frac{\epsilon}{2}
    \]
    である。

    また，
    \[
        \frac{2|\alpha|}{|\beta|^2}|b_n-\beta|
        <
        \frac{2|\alpha|}{|\beta|^2}
        \cdot
        \frac{|\beta|^2\epsilon}{4(|\alpha|+1)}
    \]
    であるから，
    \[
        \frac{2|\alpha|}{|\beta|^2}|b_n-\beta|
        <
        \frac{|\alpha|}{|\alpha|+1}\cdot\frac{\epsilon}{2}
        \leq
        \frac{\epsilon}{2}
    \]
    となる。

    よって，$n\geq N$ ならば
    \[
        \left|
            \frac{a_n}{b_n}-\frac{\alpha}{\beta}
        \right|
        \leq
        \frac{2}{|\beta|}|a_n-\alpha|
        +
        \frac{2|\alpha|}{|\beta|^2}|b_n-\beta|
        <
        \frac{\epsilon}{2}+\frac{\epsilon}{2}
        =
        \epsilon
    \]
    である。

    したがって，任意の $\epsilon>0$ に対して，ある自然数 $N$ が存在して，
    $n\geq N$ ならば
    \[
        \left|
            \frac{a_n}{b_n}-\frac{\alpha}{\beta}
        \right|<\epsilon
    \]
    が成り立つ。

    ゆえに
    \[
        \lim_{n\to\infty}\frac{a_n}{b_n}
        =
        \frac{\alpha}{\beta}
    \]
    である。
\end{proof}